\Subhead{2.1.4}{Электромагнитные волны в среде с проводимостью}

При соглашении $\e^{-\ii\omega t}$ ток $\bm j=\sigma\bm E$ включается в комплексную проницаемость $\tilde\varepsilon=\varepsilon+\ii\sigma/\omega$, где $\varepsilon$ — абсолютная диэлектрическая проницаемость. Из $\nabla\times\bm E=\ii\omega\mu\bm H$ и $\nabla\times\bm H=-\ii\omega\tilde\varepsilon\bm E$ в однородной изотропной среде следуют
\begin{equation}
 k^2=\omega^2\mu\tilde\varepsilon,\quad
 \bm H=\frac{\bm k\times\bm E}{\omega\mu},\quad
 Z=\sqrt{\frac\mu{\tilde\varepsilon}},\quad
 \overline{\bm S}=\tfrac12\operatorname{Re}(\bm E\times\bm H^*).
\end{equation}
Ветвь корня выбирается так, чтобы поле затухало по ходу передачи энергии в пассивную среду. При $\sigma\gg\omega\varepsilon$
\begin{equation}
 k\simeq\frac{1+\ii}{\delta},\qquad
 \delta=\sqrt{\frac2{\omega\mu\sigma}}.
\end{equation}
Глубина скин-слоя уменьшается как $\omega^{-1/2}$. Средние джоулевы потери $\overline q=\sigma|E|^2/2$ согласованы с убыванием потока. Нельзя пользоваться вещественным волновым сопротивлением диэлектрика в хорошем проводнике: фазы $E$ и $H$ различаются.
